% =====================================================================
%  Why Graph-Structured Systems Engaged by Finite Agents
%  Must Admit Operational Structure: A Conditional Derivation
%
%  Author: Kundai Farai Sachikonye
%          Technical University of Munich
% =====================================================================

\documentclass[11pt,a4paper]{article}

\usepackage[utf8]{inputenc}
\usepackage[T1]{fontenc}
\usepackage{lmodern}
\usepackage[a4paper,margin=2.4cm]{geometry}
\usepackage{amsmath,amssymb,amsthm,mathtools}
\usepackage{bm}
\usepackage{mathrsfs}
\usepackage{enumitem}
\usepackage{booktabs}
\usepackage{microtype}
\usepackage[round,sort&compress]{natbib}
\usepackage[colorlinks=true,linkcolor=blue!45!black,
            citecolor=blue!45!black,urlcolor=blue!45!black]{hyperref}
\usepackage{cleveref}

% ---- Theorem environments ----
\newtheoremstyle{thmstyle}{\topsep}{\topsep}{\itshape}{}{\bfseries}{.}{.5em}{}
\theoremstyle{thmstyle}
\newtheorem{theorem}{Theorem}[section]
\newtheorem{lemma}[theorem]{Lemma}
\newtheorem{proposition}[theorem]{Proposition}
\newtheorem{corollary}[theorem]{Corollary}
\theoremstyle{definition}
\newtheorem{definition}[theorem]{Definition}
\newtheorem{axiom}{Premise}
\newtheorem{principle}[theorem]{Principle}
\theoremstyle{remark}
\newtheorem{remark}[theorem]{Remark}

\crefname{axiom}{Premise}{Premises}
\Crefname{axiom}{Premise}{Premises}

% ---- Notation ----
\newcommand{\R}{\mathbb{R}}
\newcommand{\N}{\mathbb{N}}
\newcommand{\Z}{\mathbb{Z}}
\newcommand{\Zp}{\mathbb{Z}_{\geq 0}}
\newcommand{\Graph}{\mathcal{G}}
\newcommand{\Vrt}{\mathsf{V}}
\newcommand{\Edg}{\mathsf{E}}
\newcommand{\Agent}{\mathcal{A}}
\newcommand{\Op}{\mathsf{Op}}
\newcommand{\Res}{\rho}
\newcommand{\floor}{\beta}
\newcommand{\OpStr}{\mathfrak{O}}

\title{\textbf{Why Graph-Structured Systems Engaged by Finite Agents\\
Must Admit Operational Structure}\\[4pt]
\large A Conditional Derivation}

\author{%
Kundai Farai Sachikonye\\
\small Technical University of Munich\\
\small \texttt{kundai.sachikonye@wzw.tum.de}}

\date{\today}

% =====================================================================
\begin{document}
\maketitle

\begin{abstract}
\noindent
We give a conditional derivation: granted two premises---that the
engaging agents are finite, and that the engaged system is a finite
graph---we show that the system necessarily admits an
\emph{operational structure} in a precise sense. The argument is
constructive. We define a finite graph-structured system and a finite
agent, derive that any distinguishability relation an agent can
realise has a strictly positive floor on edge measure, prove that
this floor forces a residue on every agent--graph interaction, and
show that the closure of these residues under composition is a monoid
with identity and a partial order. This monoid acting on the graph is
what we call the system's operational structure. The structure is
not chosen; it is the unique closure of the agent's realisable
interactions under the constraints imposed by finiteness on both
sides. We classify operational structures into three canonical
types---algebraic, dynamical, computational---and show every
realisable structure is a combination of these. The derivation makes
no metaphysical commitments and applies wherever the two premises
hold; the framework is silent outside that scope. We close by stating
the narrow falsifiability of the main theorem: it can be refuted only
by exhibiting a finite agent realising distinguishability with zero
floor on a finite graph, which we prove is structurally impossible.
The paper is self-contained; no result of the author is cited; the
literature references are to standard mathematical lemmas.
\end{abstract}

\medskip
\noindent\textbf{Keywords:} finite agents; finite graphs; operational
structure; distinguishability floor; residue propagation; monoid
closure; conditional derivation.

\tableofcontents

% =====================================================================
\section{Introduction}
\label{sec:intro}
% =====================================================================

\subsection{The conditional question}

The question this paper addresses is short:
\begin{quote}\itshape
On the condition that the engaging agents are finite and the engaged
system is a finite graph, what structure of operations must the
system admit?
\end{quote}
We answer: a specific monoidal operational structure, derivable from
the two premises alone, with no further hypotheses.

The result is conditional throughout. We do not assert that all
agents are finite, or that all systems are graphs, or that systems
are fundamentally relational. We assert only that \emph{if} both
conditions hold, the operational structure of
\cref{thm:opstructure} is forced. The framework is silent outside
that scope: where either premise fails, the derivation does not
apply and makes no prediction.

\subsection{Why this matters}

Engineering practice routinely treats operational structures---
function composition, command pipelines, state-machine transitions,
query algebras---as design choices to be selected for convenience.
The argument here shows that for the dominant case (finite agents
engaging finite graphs, which covers essentially all of practical
software, embedded systems, sensor networks, knowledge graphs, and
finite databases) the operational structure is \emph{not} a design
choice but a derivable necessity. The class of admissible structures
is forced; only the labelling within the class is free.

This converts a large family of engineering decisions from
\emph{which structure to invent} into \emph{which structure has the
two premises already determined}. The latter is a discovery task
with a definite answer; the former is a design task with no
objective ground.

\subsection{Method}

The derivation is constructive. We fix the two premises in
\cref{sec:premises}, derive three consequences in
\cref{sec:consequences} (distinguishability floor, residue
propagation, closure under composition), prove the main theorem in
\cref{sec:main}, classify the resulting structures in
\cref{sec:classification}, and locate the falsifiability of the
argument in \cref{sec:falsifiability}. Every step is a finite
combinatorial or measure-theoretic claim with a self-contained
proof; no result of the author is cited.

\subsection{Relation to existing frameworks}

The argument touches finite graph theory \citep{diestel2017,
bollobas1998}, finite-state automata \citep{hopcroft2007}, measure
theory of finite partitions \citep{halmos1950}, monoid theory in
algebraic structures \citep{howie1995}, and order theory on finite
lattices \citep{daveypriestley2002}. We use only the textbook
content of each field. The novel contribution is the closure
argument (\cref{thm:opstructure}) that forces a particular monoidal
structure from the conjunction of the two premises; the constituent
pieces are standard, but their conjunction has not been packaged as
a forced structural result, to the author's knowledge.

% =====================================================================
\section{Premises}
\label{sec:premises}
% =====================================================================

We fix two premises. Both are operational conditions on the agent
and the system; neither is a metaphysical hypothesis about the
world.

\subsection{Finite graph-structured system}

\begin{definition}[Finite graph-structured system]
\label{def:graph}
A \emph{finite graph-structured system} is a triple $\Graph =
(\Vrt, \Edg, w)$ where $\Vrt$ is a finite non-empty set of
\emph{vertices}, $\Edg \subseteq \Vrt \times \Vrt$ is a finite set
of \emph{edges}, and $w: \Edg \to \R_{\geq 0}$ is an \emph{edge
measure} assigning to each edge a non-negative real weight. We do
not require edges to be undirected or distinct; multigraphs are
admitted.
\end{definition}

\begin{axiom}[Finite graph hypothesis]
\label{ax:graph}
The engaged system $\Graph = (\Vrt, \Edg, w)$ is a finite
graph-structured system: $|\Vrt| < \infty$, $|\Edg| < \infty$, and
$w$ takes finite values throughout.
\end{axiom}

The graph need not be connected, simple, or planar. It need not be
metric, in the sense that $w$ need not satisfy the triangle
inequality. We require only that it has a finite number of
vertices, a finite number of edges, and a finite-valued non-negative
measure on its edges.

\begin{remark}[Why graphs]
\label{rem:why_graphs}
We pick graphs as the primitive structure because they are the
minimal combinatorial encoding of relational data: a set of entities
together with the relations between them. Any system describable as
``entities standing in relations'' is encodable as a graph, possibly
a multigraph with labelled edges. Trees, finite-state machines,
knowledge bases, relational databases, dependency graphs, and finite
hypergraphs all reduce to or extend this minimal form. The choice
restricts neither the scope of application nor the generality of the
result.
\end{remark}

\subsection{Finite agent}

\begin{definition}[Distinguishability relation]
\label{def:distrel}
A \emph{distinguishability relation} on $\Graph$ is a binary
relation $\sim \subseteq \Vrt \times \Vrt$ such that for each pair
$(u, v)$, the agent has either established $u \sim v$ (the pair is
treated as a single entity), or $u \not\sim v$ (the pair is treated
as distinct entities), but not both.
\end{definition}

\begin{definition}[Finite agent]
\label{def:agent}
A \emph{finite agent} $\Agent$ on $\Graph$ is specified by:
\begin{enumerate}[label=(F\arabic*),leftmargin=2.6em,nosep]
\item\label{F:fin_state} A finite internal state space $S_\Agent$
with $|S_\Agent| < \infty$, encoding all distinctions the agent can
hold at once.
\item\label{F:fin_descr} A finite descriptive scheme: each
distinguishability decision is recorded by an injection from the
$\sim$-classes into $S_\Agent$, so the number of distinct
equivalence classes the agent can simultaneously hold is at most
$|S_\Agent|$.
\item\label{F:fin_cost} A finite cost budget: each act of
distinguishing a pair $(u,v)$ as $u \not\sim v$ has a cost
$c(u,v) > 0$ bounded below by a positive constant $c_0$, and the
agent's total expenditure over its lifetime is bounded by a finite
budget $C_{\Agent} < \infty$.
\end{enumerate}
\end{definition}

\begin{axiom}[Finite agent hypothesis]
\label{ax:agent}
Any agent engaging $\Graph$ is a finite agent in the sense of
\cref{def:agent}: it has finite internal state, finite descriptive
scheme, and finite cost budget.
\end{axiom}

We make no further assumptions about the agent's architecture: it
may be a human, a machine, a sensor network, a software process, a
finite collective. The three finiteness conditions of
\cref{def:agent} are the substantive content of \cref{ax:agent}.

\subsection{What the premises do not say}

\Cref{ax:graph,ax:agent} are sharp where they need to be sharp and
silent elsewhere. We list three things the premises explicitly do
not say.

\begin{enumerate}[nosep,leftmargin=*]
\item They do not say the world is finite. Reality at large is
permitted to be infinite; the premises restrict only the system the
agent engages and the agent itself, both of which are finite.
\item They do not say the graph encodes everything about the system.
A graph is a \emph{model} of the system the agent uses; what the
system is in itself, outside the agent's modelling, is not
addressed.
\item They do not say agents must be physical. The argument applies
to any agent satisfying~\ref{F:fin_state}--\ref{F:fin_cost},
whatever its substrate.
\end{enumerate}

\subsection{Scope}

The conjunction of the two premises covers the dominant case of
engineering practice: software systems, embedded controllers,
sensor networks, knowledge bases, finite databases, finite-state
protocols, and finite communication graphs. It does not cover
putative infinite agents (perfect oracles, unbounded computation) or
putative infinite systems (continuous fields engaged directly,
without finite discretisation). Where finite discretisation is
imposed---which is the case in all realised engineering
systems---the premises apply.

% =====================================================================
\section{Consequences of the Two Premises}
\label{sec:consequences}
% =====================================================================

We derive three consequences of the conjunction of \cref{ax:graph}
and \cref{ax:agent}. Each is forced by the premises alone; no
additional hypothesis is invoked.

\subsection{The distinguishability floor}

\begin{theorem}[Floor on edge measure for realisable
distinguishability]
\label{thm:floor}
Let $\Graph$ satisfy \cref{ax:graph} and $\Agent$ satisfy
\cref{ax:agent}. Then there exists a constant $\floor > 0$,
depending on $\Agent$ and $\Graph$ but independent of any individual
pair, such that for every pair $(u, v) \in \Vrt \times \Vrt$ that
$\Agent$ can realise as $u \not\sim v$, the path through the graph
that $\Agent$ uses to distinguish $u$ from $v$ contains an edge of
measure $\geq \floor$.
\end{theorem}

\begin{proof}
By~\ref{F:fin_cost}, every act of distinguishing has cost
$c(u,v) \geq c_0 > 0$. Suppose, for contradiction, that there is no
positive floor on the edge measures used: then for every
$\epsilon > 0$, $\Agent$ can realise a distinction $u \not\sim v$
using only edges $e$ with $w(e) < \epsilon$.

Consider the set of edge measures that appear on agent-realised
distinguishability paths. By~\ref{F:fin_descr}, the number of
distinguishability acts the agent can record is bounded by
$|S_\Agent|^2$, a finite number. Each such act uses finitely many
edges, so the union of all edge measures appearing on agent-realised
paths is a finite subset of $\R_{\geq 0}$. Null-measure edges do
not contribute to a distinction (a null edge connects vertices
that share all of the agent's path-measure, so $u \not\sim v$
across a null edge is not realisable; the agent's descriptive
scheme records no measure differential). Hence the set of
edge-measure contributions to realised distinctions is a finite
subset of $\R_{> 0}$.

The infimum of a finite non-empty set of strictly positive numbers
is itself strictly positive. Set
$\floor := \min \{w(e) : w(e) > 0,\,e \text{ appears on some
agent-realised distinguishability path}\}$. Then $\floor > 0$ and
every realised distinguishability path contains an edge of measure
$\geq \floor$, contradicting the assumption. $\qed$
\end{proof}

\begin{remark}[Floor is a finiteness floor]
\label{rem:floor_is_finiteness}
The floor $\floor$ is not a metaphysical lower bound on what edges
``can be.'' It is the lower bound forced by the finiteness of
$\Edg$ and the requirement that distinguishability be realisable.
With a finite edge set having positive-measure edges, the minimum
positive measure is itself positive. This is a combinatorial
observation, not a substantive claim about the structure of
reality.
\end{remark}

\subsection{Residue: every distinguishability act deposits content}

\begin{definition}[Residue]
\label{def:residue}
The \emph{residue} $\Res(u,v)$ of a distinguishability act
$u \not\sim v$ realised by $\Agent$ is the sum of edge measures
along the path the agent uses:
\[
  \Res(u,v) = \sum_{e \in \mathrm{path}_\Agent(u,v)} w(e).
\]
By \cref{thm:floor}, $\Res(u,v) \geq \floor > 0$ for every
realisable distinguishability act.
\end{definition}

\begin{theorem}[Residues are additive under composition]
\label{thm:residue_additive}
Let $u \not\sim v$ and $v \not\sim w$ be two distinguishability
acts realised by $\Agent$. The composed act $u \not\sim w$ (if
realisable) has residue:
\[
  \Res(u, w) \leq \Res(u, v) + \Res(v, w).
\]
Equality holds when the path the agent uses for the composed act is
the concatenation of the individual paths.
\end{theorem}

\begin{proof}
The residue of the composed act is the sum of edge measures along
the agent's path from $u$ to $w$. By the triangle inequality for
sums of non-negative numbers, the shortest such path has total
measure no greater than the sum of any concatenated route. If the
agent uses the concatenation $u \to v \to w$, equality holds;
otherwise the agent has found a shorter path and the inequality is
strict. $\qed$
\end{proof}

\begin{corollary}[Residue accumulates and does not refund]
\label{cor:res_accumulate}
For any sequence of distinguishability acts
$u_0 \not\sim u_1 \not\sim u_2 \not\sim \cdots \not\sim u_n$, the
total residue deposited satisfies:
\[
  \sum_{i=0}^{n-1} \Res(u_i, u_{i+1}) \geq n\floor > 0.
\]
The cumulative residue strictly grows with the number of acts and
is bounded below by $n\floor$. No act reduces the cumulative total;
the sum is monotone non-decreasing in $n$.
\end{corollary}

\begin{proof}
By \cref{thm:floor}, each $\Res(u_i, u_{i+1}) \geq \floor$. Summing
$n$ non-negative terms each $\geq \floor$ gives $\geq n\floor$.
Monotonicity in $n$ is immediate from the non-negativity of each
term. $\qed$
\end{proof}

\subsection{Closure: distinguishability acts compose}

\begin{definition}[Composition]
\label{def:compose}
Let $\mathcal{D}_\Agent$ be the set of distinguishability acts
$\Agent$ can realise on $\Graph$, each represented as an ordered
pair $(u, v)$ with $u \not\sim v$. \emph{Composition} of two acts
$(u, v)$ and $(v, w)$ is the act $(u, w)$, provided $\Agent$ can
realise it.
\end{definition}

\begin{lemma}[Composition is associative]
\label{lem:assoc}
Composition of distinguishability acts is associative: for any
realisable triple $(u, v)$, $(v, w)$, $(w, x)$,
\[
  ((u,v)\circ(v,w))\circ(w,x) = (u,v)\circ((v,w)\circ(w,x)).
\]
\end{lemma}

\begin{proof}
Both sides are the act $(u, x)$, identified by the same source
vertex $u$ and the same target vertex $x$. As an ordered pair, the
act is uniquely determined by its endpoints, so the two
compositions produce the same act. $\qed$
\end{proof}

\begin{lemma}[Identity exists]
\label{lem:identity}
The \emph{trivial act} $\mathrm{id} = (u, u)$ for any $u \in \Vrt$
satisfies $(u, v) \circ \mathrm{id} = (u, v)$ and
$\mathrm{id} \circ (u, v) = (u, v)$ for every realisable act
$(u, v) \in \mathcal{D}_\Agent$, where the identity is taken with
matched endpoint.
\end{lemma}

\begin{proof}
The trivial act has residue $\Res(u, u) = 0$ (no edge traversed)
and acts as identity on either side by composition. The agent need
not realise the trivial act to perform a distinguishability, but it
can include it costlessly in any composition. $\qed$
\end{proof}

\begin{theorem}[Closure under composition is a monoid]
\label{thm:monoid}
The set $\mathcal{D}_\Agent \cup \{\mathrm{id}\}$ equipped with the
composition $\circ$ forms a monoid: associative, with identity.
\end{theorem}

\begin{proof}
\cref{lem:assoc} establishes associativity; \cref{lem:identity}
establishes the identity. Closure of the underlying set under
$\circ$ holds whenever the composed act is realisable; by extending
the set to its closure under composition (adding realisable but
not directly performed acts), we obtain a monoid in the strict
sense. $\qed$
\end{proof}

% =====================================================================
\section{The Operational Structure Theorem}
\label{sec:main}
% =====================================================================

We now state and prove the main result: a finite agent engaging a
finite graph cannot avoid admitting an operational structure in the
following precise sense.

\subsection{Operational structure}

\begin{definition}[Operational structure]
\label{def:opstr}
An \emph{operational structure} on a finite graph $\Graph$ relative
to a finite agent $\Agent$ is a quadruple
$\OpStr_\Agent = (\Op, \circ, \mathrm{id}, \leq)$ where:
\begin{enumerate}[label=(O\arabic*),leftmargin=2.6em,nosep]
\item\label{O:set} $\Op$ is a finite non-empty set of
\emph{operations} the agent can apply to vertices or pairs of
vertices of $\Graph$.
\item\label{O:monoid} $\circ: \Op \times \Op \to \Op$ is an
associative composition with identity $\mathrm{id} \in \Op$.
\item\label{O:order} $\leq$ is a partial order on $\Op$ such that
$f \leq g$ iff $\Res(f) \leq \Res(g)$, where $\Res$ is the residue
deposited by the operation.
\item\label{O:floor} For every non-identity $f \in \Op$,
$\Res(f) \geq \floor > 0$.
\end{enumerate}
\end{definition}

\begin{theorem}[Operational Structure Theorem]
\label{thm:opstructure}
Under \cref{ax:graph} and \cref{ax:agent}, every distinguishability
relation $\sim$ that the agent realises on $\Graph$ induces a
unique operational structure $\OpStr_\Agent$. The structure is
forced by the conjunction of the two premises and cannot be avoided
by the agent's choice of method.
\end{theorem}

\begin{proof}
We construct $\OpStr_\Agent$ from the agent's realised
distinguishability acts and verify the four clauses.

\emph{Set of operations~\ref{O:set}.} Let $\Op = \mathcal{D}_\Agent
\cup \{\mathrm{id}\}$. By \cref{ax:agent},
$|\mathcal{D}_\Agent|$ is bounded by $|S_\Agent|^2 < \infty$.
Hence $\Op$ is finite.

\emph{Monoid structure~\ref{O:monoid}.} By \cref{thm:monoid},
$\Op$ under composition is a monoid with identity.

\emph{Partial order~\ref{O:order}.} Define $f \leq g$ iff
$\Res(f) \leq \Res(g)$. Reflexivity, transitivity, and antisymmetry
(up to residue-equivalence) follow from the same properties of
$\leq$ on $\R_{\geq 0}$.

\emph{Floor~\ref{O:floor}.} By \cref{thm:floor}, every realised
distinguishability act has residue $\geq \floor > 0$. The identity
has residue $0$, so the floor applies only to non-identity
operations, as stated.

\emph{Uniqueness.} The structure is determined by (i) the agent's
finite descriptive scheme (which fixes $\Op$ up to relabelling),
(ii) the composition law on the graph (which is determined by the
graph and the agent's path-choices, both finite), (iii) the
residue function (determined by $w$), and (iv) the floor (the
minimum positive edge measure, determined by $\Graph$). Hence the
structure is unique given $(\Agent, \Graph)$.

\emph{Forcing.} The four clauses follow from the premises alone,
with no further hypothesis. If an agent attempts to deny one of
them, the denial reduces to denying one of the premises:
rejecting~\ref{O:set} requires $|\Op| = \infty$ (denial of
\ref{F:fin_descr}); rejecting~\ref{O:monoid} requires
non-associative composition (denial of the agent's recording of the
composed result, again \ref{F:fin_descr});
rejecting~\ref{O:order} requires negative or unordered residue
(denial of \cref{def:graph}, $w \geq 0$); rejecting~\ref{O:floor}
requires zero-residue distinguishability acts (denial of
\cref{thm:floor}, hence of finiteness). The structure is forced.
$\qed$
\end{proof}

\subsection{Consequences of the theorem}

\begin{corollary}[Operations are not chosen, only labelled]
\label{cor:not_chosen}
The set $\Op$, the composition $\circ$, the order $\leq$, and the
floor $\floor$ are determined by $(\Agent, \Graph)$. Choosing an
operational structure for a finite agent on a finite graph is
choosing labels for elements of a structure that already exists; it
is not creating the structure.
\end{corollary}

\begin{proof}
Immediate from the uniqueness clause of \cref{thm:opstructure}.
The relabelling freedom is the choice of names for the operations;
the structure underneath is fixed. $\qed$
\end{proof}

\begin{corollary}[Finite-system engineering is a discovery task]
\label{cor:discovery}
Designing operations for a finite system engaged by finite agents
is the task of discovering the structure that $(\Agent, \Graph)$
determines, not of inventing one. Two designs disagreeing about
which operations to admit are disagreeing about $(\Agent, \Graph)$,
not about taste.
\end{corollary}

\begin{proof}
By \cref{cor:not_chosen}, the structure is determined; differences
between designs reduce to differences in the underlying agent--graph
pair. $\qed$
\end{proof}

\begin{corollary}[The operational structure has no minimum below the
floor]
\label{cor:no_below}
There is no operation in $\OpStr_\Agent$ with residue strictly less
than $\floor$, other than the identity. The floor is reachable
(operations of residue exactly $\floor$ may exist, corresponding to
single-edge distinctions along the lightest positive edge); it is
not crossable downward.
\end{corollary}

\begin{proof}
By~\ref{O:floor}, all non-identity operations have residue
$\geq \floor$. Operations of residue exactly $\floor$ exist iff
there is a single edge of weight $\floor$ that realises a
distinction; this is true whenever $\floor$ is attained by some
edge in $\Graph$, which by \cref{thm:floor} is the case. $\qed$
\end{proof}

% =====================================================================
\section{Classification of Operational Structures}
\label{sec:classification}
% =====================================================================

We classify the operational structures arising under the two
premises. Three canonical types exhaust the possibilities; every
realisable $\OpStr_\Agent$ is a combination of them.

\subsection{Algebraic operations}

\begin{definition}[Algebraic operation]
\label{def:alg}
An operation $f \in \Op$ is \emph{algebraic} if it acts on the
vertex set independently of the order in which its inputs are
presented, and its composition with other operations preserves
residue additivity exactly (equality in
\cref{thm:residue_additive}).
\end{definition}

Examples: equivalence-class formation, set union under a
distinguishability relation, vertex relabelling under a graph
automorphism. Algebraic operations are static: they re-express the
graph without exploiting its dynamics.

\subsection{Dynamical operations}

\begin{definition}[Dynamical operation]
\label{def:dyn}
An operation $f \in \Op$ is \emph{dynamical} if it changes the
graph itself: adding vertices, adding or removing edges, or
modifying $w$. A dynamical operation transitions
$\Graph_t \to \Graph_{t+1}$ in discrete steps indexed by
$t \in \Zp$.
\end{definition}

Examples: graph edits, state-machine transitions, edge-weight
updates. Dynamical operations exploit the temporal structure of
the agent--graph interaction.

\begin{lemma}[Dynamical operations have strictly positive residue]
\label{lem:dyn_residue}
Every realisable dynamical operation has residue $\geq \floor$.
\end{lemma}

\begin{proof}
Adding or modifying an edge of weight $< \floor$ would create an
edge that cannot contribute to any distinguishability act
(\cref{thm:floor}); the operation is vacuous. To be non-vacuous,
the operation must touch at least one edge of weight $\geq \floor$,
hence deposit residue $\geq \floor$. $\qed$
\end{proof}

\subsection{Computational operations}

\begin{definition}[Computational operation]
\label{def:comp}
An operation $f \in \Op$ is \emph{computational} if it produces a
binary output (yes/no, equal/distinct, present/absent) by traversing
a finite portion of $\Graph$ in a sequence of edge accesses.
\end{definition}

Examples: membership queries, reachability queries, path-existence
checks, equivalence tests.

\begin{lemma}[Computational operations decompose into elementary
acts]
\label{lem:comp_decompose}
Every computational operation decomposes into a finite sequence of
elementary distinguishability acts, each of residue $\geq \floor$.
The total residue is the sum of residues along the sequence.
\end{lemma}

\begin{proof}
A computational operation accesses edges in a finite sequence (by
finiteness of the agent's procedure). Each edge access either
contributes to a distinguishability act or returns the trivial act
$\mathrm{id}$. The non-trivial accesses each have residue
$\geq \floor$; the total is the sum. $\qed$
\end{proof}

\subsection{Closure under combination}

\begin{theorem}[Three-type completeness]
\label{thm:three_types}
Every realisable operation in $\OpStr_\Agent$ is a combination of
algebraic, dynamical, and computational operations. There are no
realisable operations of a fourth type.
\end{theorem}

\begin{proof}
Let $f \in \Op$ be a realisable operation. By the agent's finite
descriptive scheme, $f$ is recorded as an effect on the agent's
internal state. The effect is one of:
\begin{enumerate}[nosep,leftmargin=*]
\item A change in equivalence-class structure without graph
modification: algebraic.
\item A change in the graph itself: dynamical.
\item A query producing a binary output: computational.
\end{enumerate}
Any composite operation decomposes into these three types in some
order. There is no fourth category because the agent's effect on
the graph is exhausted by re-classifying, editing, or querying;
any other ``effect'' would have to leave both the graph and the
equivalence structure untouched while producing no output, which by
\ref{F:fin_state}--\ref{F:fin_cost} is the trivial act
$\mathrm{id}$. $\qed$
\end{proof}

\begin{remark}[Composition across types]
\label{rem:cross_type}
Operations of different types compose in $\Op$ without restriction:
a query followed by an edit followed by a relabelling is a single
composite operation. The classification labels the type of each
component; the monoid structure operates on the composites.
\end{remark}

% =====================================================================
\section{The Forced Structure in Detail}
\label{sec:detail}
% =====================================================================

We restate the main result as a structural recipe, making explicit
what $\OpStr_\Agent$ looks like once $(\Agent, \Graph)$ is fixed.

\subsection{The structure recipe}

Given $(\Agent, \Graph)$ satisfying \cref{ax:agent} and
\cref{ax:graph}:
\begin{enumerate}[label=(R\arabic*),leftmargin=2.6em,nosep]
\item\label{R:floor} Compute $\floor = \min \{w(e) : e \in \Edg,\,
w(e) > 0\}$. This is the operational floor of the structure.
\item\label{R:ops} Enumerate the agent's realisable
distinguishability acts on $\Graph$; close under composition to
obtain $\Op$.
\item\label{R:residue} For each $f \in \Op$, compute the residue
$\Res(f)$ as the path-measure sum.
\item\label{R:order} Order $\Op$ by residue: $f \leq g$ iff
$\Res(f) \leq \Res(g)$.
\item\label{R:classify} Classify each $f \in \Op$ as algebraic,
dynamical, or computational by inspection of its effect.
\end{enumerate}

\begin{proposition}[The recipe is well-defined and terminates]
\label{prop:recipe}
The recipe \ref{R:floor}--\ref{R:classify} is well-defined and
terminates in finite time for any $(\Agent, \Graph)$ satisfying the
two premises.
\end{proposition}

\begin{proof}
\ref{R:floor} is a minimum over a finite non-empty set of positive
reals: well-defined and computable in $O(|\Edg|)$ time.
\ref{R:ops} is the closure of a finite generating set under a
finite binary operation: terminates because
$|\Op| \leq |S_\Agent|^2$ is finite.
\ref{R:residue} sums finitely many non-negative reals per
operation: well-defined and computable in $O(|\Vrt|)$ per operation.
\ref{R:order} is a comparison on a finite totally ordered set
($\R_{\geq 0}$): well-defined.
\ref{R:classify} is a case distinction over three cases per
operation: well-defined.
Each step terminates; the recipe terminates. $\qed$
\end{proof}

\subsection{What the structure looks like}

\begin{proposition}[Structural shape]
\label{prop:shape}
The operational structure $\OpStr_\Agent$ has the following shape:
\begin{enumerate}[nosep,leftmargin=*]
\item It is a finite monoid (\cref{thm:monoid}).
\item Its non-identity elements have residues bounded below by
$\floor > 0$ (\cref{thm:floor}).
\item It is partially ordered by residue, with the identity as the
unique minimum (residue $0$) and the maximum being the operation of
highest residue (the most expensive realised act).
\item Its operations decompose into three canonical types
(\cref{thm:three_types}).
\end{enumerate}
\end{proposition}

\begin{proof}
Each clause is a restatement of an earlier result. The shape is the
conjunction of those results. $\qed$
\end{proof}

\subsection{Why the structure is operational rather than only
combinatorial}

The structure $\OpStr_\Agent$ is called \emph{operational} because
each element represents an act the agent can perform, with a
quantifiable cost (residue) and a quantifiable order (residue
comparison). It differs from a purely combinatorial structure---a
finite monoid considered abstractly---in two ways:
\begin{enumerate}[nosep,leftmargin=*]
\item Each element is associated with a concrete act on the graph,
not merely an abstract symbol.
\item Composition is realised, not merely defined: composing two
acts produces a third act the agent can perform, with predictable
residue deposited on the graph.
\end{enumerate}
The operational character is essential to the result: we are not
classifying the agent's grammar but its realisable acts on a
concrete finite system. The two coincide only because the premises
force the agent's grammar to track its acts; in richer settings
the two diverge.

% =====================================================================
\section{Falsifiability}
\label{sec:falsifiability}
% =====================================================================

\subsection{Where the argument can fail}

The main theorem (\cref{thm:opstructure}) is a conditional claim
under \cref{ax:graph} and \cref{ax:agent}. It can be falsified only
by:
\begin{enumerate}[nosep,leftmargin=*]
\item Showing the two premises jointly imply a contradiction (they
do not, by inspection: a finite graph and a finite agent obviously
coexist; we are reading this on one).
\item Showing some step in the derivation does not follow from the
premises (a mathematical refutation of one of the proved lemmas).
\item Showing the premises do not apply to a particular case (which
is honest: the framework declines to address that case).
\end{enumerate}

The substantive falsification target is (2): exhibiting a finite
agent realising distinguishability with zero floor on a finite
graph. We close this case in the next subsection.

\subsection{Zero-floor distinguishability is structurally impossible}

\begin{theorem}[Zero-floor impossibility]
\label{thm:zero_floor}
Under \cref{ax:graph} and \cref{ax:agent}, there is no agent--graph
pair $(\Agent, \Graph)$ admitting a realisable distinguishability
act of residue $< \floor$, where $\floor$ is the operational floor
defined by the pair.
\end{theorem}

\begin{proof}
Suppose for contradiction such a pair exists, with an act
$u \not\sim v$ of residue $\Res(u, v) < \floor$. Since the residue
is the sum of edge measures along the agent's path, and each edge
measure is either $0$ (does not contribute) or $\geq \floor$ (by
the definition of $\floor$ as the minimum positive edge measure
appearing on some realised path), the sum is either $0$ or
$\geq \floor$.

If the sum is $0$, the agent has used only null edges: the two
vertices $u, v$ are connected by null edges only, hence are
indistinguishable in $\Graph$ (any path the agent could take has
zero measure, so by~\ref{F:fin_descr}, the agent cannot record a
non-trivial distinguishability between them). The act is vacuous;
it is not a realised distinguishability.

If the sum is $\geq \floor$, the residue is $\geq \floor$,
contradicting the assumption.

In either case, the assumed pair does not exist. $\qed$
\end{proof}

\begin{corollary}[The main theorem is irrefutable from within the
premises]
\label{cor:irrefutable}
\Cref{thm:opstructure} cannot be refuted by an example satisfying
\cref{ax:graph} and \cref{ax:agent}. Refutation requires either
denying one of the premises or finding an error in a constituent
lemma; no example produced under the premises is a counterexample.
\end{corollary}

\begin{proof}
Refutation by example requires a $(\Agent, \Graph)$ satisfying the
premises in which the theorem fails. By the construction of
$\OpStr_\Agent$ from $(\Agent, \Graph)$ in the theorem's proof, and
the zero-floor impossibility (\cref{thm:zero_floor}), no such
example exists. $\qed$
\end{proof}

\subsection{Honest scope}

\begin{remark}[The argument is silent outside its premises]
\label{rem:silent}
We make no claim about agents that fail \cref{ax:agent} (putative
infinite agents, oracles, perfect observers) or systems that fail
\cref{ax:graph} (continuous fields without finite discretisation,
infinite combinatorial structures). The argument has nothing to
say about those cases. This silence is intended: the framework is
deliberately narrow, sharp where the premises hold, mute elsewhere.
This is the right shape for a conditional architectural claim. It
allows the result to be definite without overreaching.
\end{remark}

% =====================================================================
\section{Implications}
\label{sec:implications}
% =====================================================================

We draw three implications of the main theorem for engineering
practice. Each is a direct consequence of \cref{thm:opstructure} or
its corollaries.

\subsection{Software architecture under the conditional}

The conjunction of \cref{ax:graph} and \cref{ax:agent} covers all
realised software systems: any program with finite memory and
finite execution is a finite agent, and any data structure it
manipulates is a finite graph. By \cref{thm:opstructure}, the
program's operational structure is forced by its data structure and
its agency: it is not chosen.

In practice this means: arguments about which operations a program
``should'' admit reduce to arguments about $(\Agent, \Graph)$.
Two architectures disagreeing about the operational interface are
disagreeing about one of: which finite data structure the system is
modelled as; which finite agent the program is. Once those are
agreed, the operations are determined.

\subsection{Database query algebras}

A finite database is a finite graph in the sense of
\cref{def:graph}: the schema gives the vertices (tables, attributes,
values), the foreign-key relations give the edges, and the
cardinalities give the edge measures. A finite query engine is a
finite agent in the sense of \cref{def:agent}: bounded internal
state, finite query plan space, finite cost budget per query.

By \cref{thm:opstructure}, the query algebra is forced. The choice
between relational algebra, calculus, and CRUD is the choice of
which subset of the forced algebra to expose; it is not the choice
of which algebra to invent. The minimal operational floor is the
smallest positive edge weight in the schema graph (typically a unit
cardinality), and every realisable query has residue $\geq$ this
floor.

\subsection{Knowledge graphs and reasoning}

A finite knowledge graph---ontology, RDF triple store, labelled
property graph---is a finite graph in the sense of \cref{def:graph}.
A finite reasoner---SPARQL engine, description-logic prover, graph
neural network---is a finite agent in the sense of \cref{def:agent}.
By \cref{thm:opstructure}, the reasoner's operational structure is
forced by the graph it queries and the resources it has.

The implication is that the apparent diversity of reasoning systems
collapses, under the conditional, to relabelling. Two reasoners
with the same $(\Agent, \Graph)$ admit the same operations up to
naming; their differences are presentational, not structural.
Cross-system calibration (translation of one query language into
another) is guaranteed to exist whenever $(\Agent, \Graph)$ is
shared.

% =====================================================================
\section{Conclusion}
\label{sec:conclusion}
% =====================================================================

\subsection{What we have established}

We have shown that under two finiteness premises---finite engaging
agent (\cref{ax:agent}), finite engaged graph-structured system
(\cref{ax:graph})---the system necessarily admits an operational
structure $\OpStr_\Agent$ that is forced by the premises and unique
given the pair $(\Agent, \Graph)$. The structure is a finite monoid
with a partial order by residue, a strictly positive floor on
non-identity residues, and a three-type classification (algebraic,
dynamical, computational) that exhausts the realisable operations.

The result is a conditional architectural claim. It does not
address what reality is in itself; it addresses what structure
follows when the two conditions hold. Where they hold---which is
the case in essentially all realised engineering systems---the
operational structure is determined.

\subsection{What we have not claimed}

We have not claimed that:
\begin{enumerate}[nosep,leftmargin=*]
\item Reality is fundamentally finite. The premises restrict the
agent and the system the agent engages, not the world.
\item All systems are graphs. The premises identify the conditions
under which the result holds; systems failing the conditions are
not addressed.
\item Operational structures are the only structures of interest.
Other structures (continuous dynamics, categorical limits, infinite
processes) may coexist with operational structure; we have not
spoken to them.
\end{enumerate}

\subsection{The narrowness of the result is the point}

A conditional result is sharp because it is narrow. The framework
established here speaks definitively about a specific class of
situations and is silent about others. This is exactly the right
shape for an architectural claim: scope is a feature, and the scope
of \cref{thm:opstructure} is the scope of finite agency engaging
finite graph-structured systems. Within that scope, the operational
structure is forced. Outside it, the framework declines to address
the case.

\subsection{What this enables}

The main practical consequence is that engineering decisions over
$(\Agent, \Graph)$ become discovery tasks rather than invention
tasks. To design a finite system's operational structure is to
identify the unique $\OpStr_\Agent$ that the agent--graph pair
determines. To compare two designs is to ask whether they realise
the same $(\Agent, \Graph)$. To extend a system is to ask how the
extension changes $(\Agent, \Graph)$ and re-derive the structure.

The framework converts what is currently treated as a matter of
architectural taste into a matter of formal derivation. Whether
this is a useful conversion in practice depends on whether the cost
of formalising $(\Agent, \Graph)$ for a given system is less than
the cost of debating its operational structure informally. We
expect this to be true at scale---for large systems with many
stakeholders and long lifetimes---and false at small scale, where
the informal debate is cheap.

\subsection{Closing}

The structure of operations on a finite graph engaged by a finite
agent is not a design problem. It is a derivation from the two
finiteness premises. The derivation produces a unique answer. The
answer is a finite monoid with a positive floor and a three-type
classification. That is what the conditional architecture says, and
that is all it says.

% =====================================================================
\begin{thebibliography}{9}

\bibitem[Bollob\'as(1998)]{bollobas1998} B\'ela Bollob\'as.
\emph{Modern Graph Theory}.
Graduate Texts in Mathematics 184, Springer, 1998.

\bibitem[Davey and Priestley(2002)]{daveypriestley2002}
B.~A. Davey and H.~A. Priestley.
\emph{Introduction to Lattices and Order}, 2nd edition.
Cambridge University Press, 2002.

\bibitem[Diestel(2017)]{diestel2017} Reinhard Diestel.
\emph{Graph Theory}, 5th edition.
Graduate Texts in Mathematics 173, Springer, 2017.

\bibitem[Halmos(1950)]{halmos1950} Paul R. Halmos.
\emph{Measure Theory}.
Van Nostrand, 1950.

\bibitem[Hopcroft, Motwani and Ullman(2007)]{hopcroft2007}
John E. Hopcroft, Rajeev Motwani, and Jeffrey D. Ullman.
\emph{Introduction to Automata Theory, Languages, and Computation},
3rd edition. Addison-Wesley, 2007.

\bibitem[Howie(1995)]{howie1995} John M. Howie.
\emph{Fundamentals of Semigroup Theory}.
London Mathematical Society Monographs 12, Oxford University Press,
1995.

\end{thebibliography}

\end{document}
